% ══════════════════════════════════════════════════════════════════════════════
%  CHAPTER 1: THE EXPANDING UNIVERSE
% ══════════════════════════════════════════════════════════════════════════════

\chapter{The Expanding Universe}

% ──────────────────────────────────────────────────────────────────────────────
\section{Hubble's Discovery}

\subsection{Physics Intuition}

\begin{intuition}
In the early twentieth century, astronomers noticed that the spectral lines of distant
galaxies were systematically shifted toward the red end of the spectrum.
Edwin Hubble's landmark 1929 observations at the Mount Wilson Observatory revealed a
clear linear relation between a galaxy's recession velocity and its distance from us.
This was the first observational evidence that the universe is not static but
expanding—every galaxy is moving away from every other, like raisins in a rising loaf
of bread. The discovery overturned centuries of assumptions about a steady, eternal
cosmos and laid the empirical foundation for modern Big Bang cosmology.
\end{intuition}

\subsection{Theory}

\begin{theory}
To quantify Hubble's discovery we infer recession velocity from the Doppler-like
redshift of a galaxy's spectral lines: $v = cz$ in the non-relativistic limit.
Physical distances are estimated via the \emph{cosmic distance ladder}—at
increasing scales: parallax, Cepheid variable stars, and Type~Ia supernovae.

The \textbf{Hubble--Lema\^itre Law} relates the recession velocity of a galaxy to its
proper distance~$d$:
\begin{equation}
  v \;=\; H_0 \, d,
  \label{eq:hubble}
\end{equation}
where $H_0 \approx 67.4 \kms\Mpc^{-1}$ is the Hubble constant (Planck~2018).
The cosmological redshift is defined via the scale factor $a(t)$:
\begin{equation}
  1 + z \;=\; \frac{a(t_0)}{a(t_e)} \;=\; \frac{1}{a(t_e)},
  \label{eq:redshift}
\end{equation}
using the convention $a(t_0)=1$ today.

The Hubble constant $H_0$ encodes the \emph{current} expansion rate.
Direct distance-ladder measurements give $H_0 \approx 73\,\mathrm{km\,s^{-1}\,Mpc^{-1}}$,
while CMB-based inference (Planck~2018) yields $H_0 \approx 67.4\,\mathrm{km\,s^{-1}\,Mpc^{-1}}$.
This $\sim5\sigma$ discrepancy—the \emph{Hubble tension}—is one of the most pressing
open problems in modern cosmology.
\end{theory}

\subsection{Question Examples}

\begin{example}{Recession Velocity from Redshift}{}
A galaxy in the Virgo Cluster is observed with a spectroscopic redshift of $z = 0.0036$.
Using the Hubble--Lema\^itre Law with $H_0 = 67.4\,\mathrm{km\,s^{-1}\,Mpc^{-1}}$,
estimate (a)~its recession velocity and (b)~its distance.

\medskip\textbf{Solution.}\quad
(a)~In the non-relativistic limit, $v = cz = (3\times10^5)(0.0036)
= 1080\,\mathrm{km\,s^{-1}}$.
(b)~From $v = H_0 d$, $d = 1080/67.4 \approx 16\,\mathrm{Mpc}$, consistent with
the known Virgo Cluster distance of $\sim16$--$17\,\mathrm{Mpc}$.
\end{example}

\subsection{Extra}

\begin{extra}
Hubble's original 1929 determination yielded $H_0 \approx 500\,\mathrm{km\,s^{-1}\,Mpc^{-1}}$—nearly
eight times the modern value—because he confused two classes of Cepheid variable stars:
classical (Population~I) Cepheids and the intrinsically dimmer W~Virginis stars.
Walter Baade's 1952 revision of the Cepheid calibration corrected the distance scale and
reconciled the estimated age of the universe with the known ages of the oldest globular clusters.
\end{extra}

% ──────────────────────────────────────────────────────────────────────────────
\section{The Cosmological Redshift}

\subsection{Physics Intuition}

\begin{intuition}
The cosmological redshift is fundamentally different from the ordinary Doppler
redshift caused by a galaxy's peculiar velocity through space.
Instead, it is \emph{space itself} that expands, stretching the wavelength of every
photon in transit proportionally to how much the universe has expanded during the photon's
journey. The longer a photon travels, the greater its observed redshift—making $z$ a
direct measure of cosmic distance and look-back time.
\end{intuition}

\subsection{Theory}

\begin{theory}
Consider a photon emitted at cosmic time $t_e$ with wavelength $\lambda_e$, received
today ($t_0$) with wavelength $\lambda_0$.
In the FLRW metric, the null geodesic condition $ds^2=0$ for radial travel implies
that the proper wavelength scales with the scale factor:
$\lambda \propto a(t)$.
Therefore
\begin{equation}
  \frac{\lambda_0}{\lambda_e} \;=\; \frac{a(t_0)}{a(t_e)} \;=\; 1 + z,
\end{equation}
reproducing Eq.~\eqref{eq:redshift}.
This derivation requires only the FLRW metric—no specific model for $a(t)$ is needed.
In the limit of small separations, $z \ll 1$, it reduces to the familiar Doppler
formula $z \approx v/c$ with $v = H_0 d$.
\end{theory}

\subsection{Question Examples}

\begin{example}{Computing Cosmic Distance}{}
A hydrogen Lyman-$\alpha$ photon (rest wavelength $\lambda_e = 121.6\,\mathrm{nm}$)
from a distant quasar is detected at $\lambda_0 = 729.6\,\mathrm{nm}$.
Find the redshift $z$ and the scale factor $a(t_e)$ at the time of emission.

\medskip\textbf{Solution.}\quad
$z = \lambda_0/\lambda_e - 1 = 729.6/121.6 - 1 = 5.0$.
From $1+z = 1/a(t_e)$, we get $a(t_e) = 1/6 \approx 0.167$:
the universe was six times smaller when this photon was emitted.
\end{example}

\subsection{Extra}

\begin{extra}
The most distant objects confirmed to date reach redshifts $z > 13$, placing them
within the first few hundred million years after the Big Bang.
JWST has revealed surprisingly luminous, massive galaxies at $z > 10$, challenging
standard semi-analytic models of early galaxy formation.
Such observations probe the \emph{epoch of reionization}—when ultraviolet radiation
from the first stars and quasars ionized the neutral hydrogen permeating the universe.
\end{extra}

\clearpage
